% ============================================================================
% Apêndice: Algoritmo do Grafo de Chamadas (Call Graph), versão corrigida
% para refletir a implementação real (call_graph_builder.py,
% definition_index.py, import_index.py, render_source_view.py).
%
% Requer no preâmbulo (se ainda não existir): \usepackage{listings}
%
% Principais correções em relação ao rascunho anterior:
%
% 1. O nome "Structural Code Summarizer" e a ideia de "sumarização" foram
%    removidos. O algoritmo real NÃO resume o código (não preserva "só
%    assinaturas e tipos", não omite "boilerplate"). Para cada nó resolvido
%    da árvore, ele lê e inclui o trecho de código-fonte LITERAL do
%    repositório clonado em disco (render_source_view.py). O único corte que
%    existe é um limite bruto de caracteres aplicado depois, na Etapa 3
%    (MAX_SOURCE_CODE_CHARS = 32.000, em evaluation/payload.py), para conter
%    o custo/atenção do juiz. É truncamento, não sumarização.
%
% 2. As "3 fases" foram reescritas para bater com o código real: (1)
%    indexação estática do repositório (índice de definições + índice de
%    imports por arquivo, não um "mapeamento de entry point"); (2)
%    construção RECURSIVA da árvore (não BFS) com uma heurística de
%    resolução de 5 passos por nome, primeira que casar vence; (3)
%    renderização do contexto via leitura literal do código-fonte de cada
%    nó resolvido, não uma "linearização" que preserva só estrutura.
%
% 3. O pseudocódigo foi reescrito para espelhar as funções reais
%    (build_call_graph, resolve_call, render_source_view) em vez de nomes
%    genéricos inventados (extract_call_graph, get_internal_function_calls,
%    format_to_tree_view).
%
% 4. A citação \cite{shi2026description} foi REMOVIDA: não há como
%    verificar se essa referência existe de fato nem se sustenta a alegação
%    original ("intenção lógica definida na camada superior da pilha"), e eu
%    não vou inventar uma citação. A justificativa para profundidade=3 foi
%    reescrita em cima do que o próprio código documenta (contenção de
%    custo + eliminação estrutural de recursão infinita), sem fonte externa.
%    Se vocês tiverem uma referência real para essa alegação, adicionem a
%    citação de volta manualmente.
%
% 5. \label{code:scs} -> \label{code:call-graph} (o nome antigo vinha do
%    acrônimo removido). Ainda não deve haver \ref para ele em outro lugar,
%    já que este apêndice parece ser rascunho novo, mas vale conferir.
% ============================================================================

\section{Algoritmo de Construção do Grafo de Chamadas e Composição do Contexto de Código}
\label{apendice:secao}

Este apêndice detalha o algoritmo desenvolvido para a construção do grafo de
chamadas (\textit{call graph}) de cada ferramenta MCP e para a composição do
trecho de \textcolor{green}{código} usado como contexto adicional no cenário de avaliação
que inclui \textcolor{green}{código} (Seção~\ref{sec:materiais}). O algoritmo é
heurístico e baseado em análise estática via tree-sitter: a partir da função
que implementa a própria ferramenta, resolve por nome as chamadas feitas
dentro dela, até três níveis de profundidade, distinguindo chamadas
resolvidas dentro do próprio repositório de chamadas externas ou
dinamicamente determinadas, conforme ilustrado na Figura~\ref{fig:codigoarvore}.

\subsection*{Procedimento de Extração}
O algoritmo opera em três fases distintas, executadas uma vez por repositório:

\begin{enumerate}
    \item \textbf{Indexação estática do repositório:} são construídos, a
    partir de todo o \textcolor{green}{código} do repositório, um índice de definições
    (nome simples e nome qualificado, \texttt{Classe.método}, de toda
    função/método encontrado, com referência ao seu corpo) e, por arquivo,
    um índice de \textit{imports} (alias local usado no arquivo
    $\rightarrow$ módulo e nome original de onde veio). Ambos são
    heurísticas de escopo léxico, não uma resolução completa de símbolos.

    \item \textbf{Construção recursiva do grafo de chamadas:} a árvore de
    cada ferramenta é construída a partir da função que a implementa
    (nível~1). Para cada chamada encontrada no corpo de uma função, uma
    heurística de resolução por nome é aplicada em ordem, valendo a
    primeira que casar: (i) receptor \texttt{self}/\texttt{this}
    $\rightarrow$ método da própria classe; (ii) mesmo nome definido no
    mesmo arquivo $\rightarrow$ função auxiliar local; (iii) receptor casa
    com um alias de \textit{import} conhecido $\rightarrow$ busca no
    módulo importado; (iv) nome único em todo o repositório
    $\rightarrow$ resolve diretamente; (v) não encontrado em lugar nenhum
    (biblioteca externa, \textit{stdlib}, ou chamada genuinamente
    dinâmica) $\rightarrow$ tratado como folha externa, sem recursão.
    Uma chamada resolvida vira um novo nó, expandido recursivamente até o
    nível~3; nós de nível~3 continuam sendo chamadas normalmente
    resolvidas, mas suas próprias chamadas internas não são mais
    exploradas. Quando mais de um candidato válido é encontrado num mesmo
    passo (nome duplicado em escopos distintos), o desempate é feito pela
    menor distância de diretório até o arquivo de origem da chamada, e o
    nó resultante permanece explicitamente sinalizado como \textit{ambíguo}
    na estrutura de dados, em vez de ser tratado como uma resolução certa.

    \item \textbf{Composição do contexto de código:} no cenário de
    avaliação que inclui \textcolor{green}{código}, a árvore é percorrida em
    pré-ordem e, para cada nó resolvido, o trecho de \textcolor{green}{código}
    correspondente (linhas de início e fim exatas) é lido diretamente do
    repositório clonado em disco e anotado com um comentário indicando o
    nível na árvore, a chamada de origem e a localização exata
    (arquivo:linhas). Nós externos/não resolvidos entram apenas como um
    comentário sinalizando a chamada não resolvida, sem trecho de código.
    Os trechos são concatenados na ordem de visita; o texto resultante é
    limitado a um número máximo de caracteres antes de ser enviado ao
    avaliador, para conter o custo e a atenção do modelo: um corte por
    tamanho, não uma sumarização do conteúdo.
\end{enumerate}

\subsection*{Pseudocódigo do Algoritmo}
A lógica de construção do grafo e de resolução de chamadas é apresentada na
Listagem~\ref{code:call-graph}. A profundidade máxima é imposta pela
própria recursão: ao atingir o nível limite, a função retorna sem procurar
novas chamadas, o que também elimina, por construção, qualquer risco de
loop em recursão direta ou mútua, sem exigir detecção de ciclo separada.

\begin{lstlisting}[language=Python, caption=Construção do grafo de chamadas e resolução por nome, label=code:call-graph]
MAX_LEVEL = 3

def build_call_graph(fn_def, level=1):
    node = Node(level, resolved=True, external=False,
                ambiguous=False, location=fn_def.location)

    if level >= MAX_LEVEL:
        return node  # no de nivel 3: resolvido, mas nao expandido

    for call_site in extract_calls(fn_def.body):
        resolved_def, ambiguous = resolve_call(call_site, fn_def)

        if resolved_def is None:
            node.calls.append(
                ExternalLeaf(level + 1, call_site.raw_text))
            continue

        child = build_call_graph(resolved_def, level + 1)
        child.ambiguous = ambiguous
        child.raw_call_text = call_site.raw_text
        node.calls.append(child)

    return node


def resolve_call(call_site, fn_def):
    name = call_site.callee_name

    # 1. receptor self/this -> metodo da propria classe
    if call_site.receiver in ("self", "this"):
        candidate = definitions.by_qualified_name.get(
            f"{fn_def.class_name}.{name}")
        if candidate:
            return candidate, False  # (definicao, ambiguous)

    # 2. mesmo nome, mesmo arquivo
    # 3. receptor casa com alias de import conhecido
    # 4. nome unico em todo o repositorio
    candidates = lookup_by_scope(name, call_site, fn_def)
    if len(candidates) == 1:
        return candidates[0], False
    if len(candidates) > 1:
        return nearest_by_directory(candidates, fn_def.file), True

    # 5. nao encontrado -> chamada externa/dinamica
    return None, False
\end{lstlisting}

\subsection*{Justificativa Técnica}
A limitação a três níveis é uma decisão de engenharia, não uma medição
prévia de onde a lógica relevante de uma ferramenta termina: ela contém o
tamanho da árvore (e, por consequência, do texto eventualmente enviado ao
avaliador) a um custo previsível, e elimina, por construção, qualquer
risco de loop infinito em recursão direta ou mútua, sem exigir detecção de
ciclo separada. A resolução de chamadas é heurística e baseada em nome, não
uma checagem de tipos ou resolução semântica completa: um \textit{trade-off}
assumido dado o escopo do projeto. Por isso, toda chamada resolvida por
desempate entre múltiplos candidatos válidos permanece marcada como
ambígua na estrutura de dados final, em vez de ser silenciosamente tratada
como uma resolução certa, preservando essa incerteza para análises
posteriores.

\section{Algoritmo de Cálculo de LOC e Complexidade Ciclomática da Tool Completa}
\label{apendice:secao-complexidade}

\textcolor{green}{Este apêndice detalha o algoritmo usado para calcular as métricas de
tamanho (LOC) e complexidade ciclomática de cada ferramenta, apresentadas na
Seção~\ref{sec:materiais}. Ele reaproveita a heurística de resolução de
chamadas por nome descrita no Apêndice~\ref{apendice:secao} (mesmos cinco
passos, mesma estrutura de índices), mas resolve um problema diferente: em
vez de uma árvore de profundidade fixa usada para compor o contexto de
código enviado ao avaliador, aqui o objetivo é determinar o conjunto
completo de funções que compõem a implementação de uma ferramenta, sem um
teto de profundidade arbitrário.}

\subsection*{Procedimento}

\begin{enumerate}
    \item \textbf{\textcolor{green}{Busca em largura sem teto de profundidade:}}
    \textcolor{green}{partindo da função que implementa a ferramenta, cada chamada
    encontrada em seu corpo é resolvida pela mesma heurística de nome do
    Apêndice~\ref{apendice:secao}. Toda função resolvida e ainda não
    visitada é adicionada ao conjunto de funções alcançáveis e também tem
    seu próprio corpo percorrido em busca de novas chamadas, repetindo-se o
    processo até não haver mais chamadas resolvidas e não visitadas. Um
    conjunto de nomes qualificados já visitados evita tanto a recontagem de
    uma mesma função alcançada por mais de um caminho (por exemplo, duas
    funções distintas que chamam um mesmo utilitário) quanto um loop
    infinito em caso de recursão direta ou mútua -- papel que, na árvore de
    três níveis do Apêndice~\ref{apendice:secao}, é desempenhado pelo
    próprio teto de profundidade.}

    \item \textbf{\textcolor{green}{Restrição adicional para chamadas com receptor:}}
    \textcolor{green}{diferentemente da árvore de três níveis, uma chamada da forma
    \texttt{objeto.metodo(...)} cujo receptor não é resolvido por nenhuma
    heurística baseada em contexto local (\texttt{self}/\texttt{this} ou
    alias de \textit{import} conhecido) não avança para o último passo da
    heurística (nome único em todo o repositório): ela é tratada como
    externa, do mesmo modo que uma chamada cujo nome não é encontrado em
    lugar nenhum. Chamadas sem receptor (\texttt{helper(...)}) continuam
    elegíveis a esse último passo normalmente. Essa restrição existe porque,
    sem teto de profundidade, um único acerto incorreto desse passo --
    aceitável como uma curiosidade isolada e visível num nó da árvore de
    três níveis -- se propagaria por todas as chamadas feitas a partir da
    função incorretamente associada, e assim sucessivamente. Um caso real,
    encontrado num repositório que empacota dezenas de servidores MCP não
    relacionados sob uma única árvore de código-fonte, tornou esse risco
    concreto: a chamada \texttt{client.p.submission.fetch(...)} (um método
    de uma biblioteca cliente externa) foi ``resolvida'' para uma função
    \texttt{fetch} de um servidor completamente não relacionado, pelo simples
    fato de ser a única função com esse nome em todo o repositório. Chamadas
    sem receptor continuam usando esse último passo porque é, na prática, a
    única forma de a heurística acompanhar um padrão de importação comum em
    JavaScript/TypeScript (\texttt{const \{ helper \} = require(...)}), não
    capturado pelo índice de \textit{imports} construído na fase de
    indexação estática.}

    \item \textbf{\textcolor{green}{Contagem por função:}} \textcolor{green}{para cada
    função do conjunto alcançável, LOC é a diferença entre sua linha final e
    inicial mais um, e a complexidade ciclomática é a contagem de McCabe
    \cite{mccabe1976complexity} (pontos de decisão mais um) sobre seu
    próprio corpo, sem descer em funções ou \textit{lambdas} aninhadas
    definidas dentro dele. Os pontos de decisão contados variam por
    linguagem, mas seguem o mesmo princípio em todas: \texttt{if}
    (e variantes de encadeamento, como \texttt{elif}/\texttt{else if}),
    laços (\texttt{for}/\texttt{while}/\texttt{do-while}), captura de
    exceção (\texttt{catch}/\texttt{except}/\texttt{rescue}), cada ramo de
    um \texttt{switch}/\texttt{match}/\texttt{when} exceto o ramo
    \textit{default}, e expressões ternárias -- nunca operadores lógicos de
    curto-circuito (\texttt{\&\&}/\texttt{||}), o que caracteriza a
    definição clássica de McCabe, e não a variante estendida usada por
    algumas ferramentas de análise estática. Os totais de LOC e de
    complexidade ciclomática da ferramenta são, então, a soma desses
    valores sobre todas as funções do conjunto alcançável.}
\end{enumerate}

\subsection*{Pseudocódigo do Algoritmo}

\begin{lstlisting}[language=Python, caption={Fecho transitivo de chamadas resolvidas, sem teto de profundidade}, label=code:reachable-closure]
def collect_reachable_definitions(start_def):
    visited = {start_def.qualified_name: start_def}
    queue = [start_def]

    while queue:
        current = queue.pop()

        for call_site in extract_calls(current.body):
            # unico ponto de diferenca em relacao a resolve_call()
            # do Apendice A: chamada com receptor nao reconhecido
            # nao cai no passo 4 (nome unico no repositorio inteiro)
            resolved_def, ambiguous = resolve_call(
                call_site, current,
                allow_repo_wide_fallback=(call_site.receiver is None))

            if resolved_def is None or ambiguous:
                continue  # externa/dinamica, ou resolucao incerta: nao conta
            if resolved_def.qualified_name in visited:
                continue  # ja alcancada por outro caminho, ou recursao

            visited[resolved_def.qualified_name] = resolved_def
            queue.append(resolved_def)

    return list(visited.values())


def loc_e_complexidade(reachable_defs):
    loc_total, cc_total = 0, 0
    for fn_def in reachable_defs:
        loc_total += fn_def.end_line - fn_def.start_line + 1
        cc_total += count_decision_points(fn_def.body) + 1  # McCabe
    return loc_total, cc_total
\end{lstlisting}

\subsection*{Justificativa Técnica}

\textcolor{green}{A ausência de um teto de profundidade aqui, em contraste direto com o
Apêndice~\ref{apendice:secao}, é deliberada: o objetivo dessa métrica é
estimar o tamanho e a complexidade reais de uma ferramenta, e um teto fixo
sub-representaria sistematicamente ferramentas cuja lógica está distribuída
por mais de três níveis de chamadas internas. A restrição adicional sobre
chamadas com receptor no passo de resolução (item 2 acima) é o
\textit{trade-off} inverso: sem teto de profundidade, a mesma heurística
best-effort baseada em nome que é inofensiva numa árvore de três níveis
passa a poder inflar uma métrica agregada por um fator arbitrariamente
grande a partir de um único acerto equivocado, o que a tornaria não
confiável para qualquer análise comparativa entre ferramentas.}
